\documentclass{my_eth_beamer}

\usepackage{my_math_eth}

\numberwithin{equation}{section}

% -----------------------------------------------------------------------------
% this document
\title{Introduction to Survey Sampling}
\newcommand{\TITLE}{Calibration}
\newcommand{\NR}{8}
\subtitle{Part \NR: \TITLE}
\author[Tobias Schoch]{Tobias Schoch}
\date{2026}

% -----------------------------------------------------------------------------
\begin{document}

\setstretch{1.1}

\section{\TITLE}\label{sec:calib}
\setcounter{section}{\NR}
\setcounter{Bsp}{19}
\setcounter{Def}{13}
\setcounter{EinSatz}{1}

% -----------------------------------------------------------------------------
\begin{frame}
    \frametitle{Overview}
    \noindent \hspace{0.2em}\textcolor{fhnwblue}{\bfseries \NR~\TITLE}
    \vspace{0.25em}
    \begin{enumerate}[\NR.1]
    \setlength{\itemsep}{4pt}
        \item Information, groups, and domains
        \item Post-stratification
        \item Calibration
    \end{enumerate}
\end{frame}

\subsection{Information, groups, and domains}

% -----------------------------------------------------------------------------
\begin{frame}
    \frametitle{\insertsectionnumber.\insertsubsectionnumber~\insertsubsection}
    \begin{itemize}
    \setlength{\itemsep}{6pt}
        \item We consider \alert{integration of information} about the
            population with the aim of:
            \vspace{0.25em}
            \begin{itemize}
            \setlength{\itemsep}{3pt}
                \item stratification or unequal-probability sampling
                    ($\Rightarrow$ efficiency gains),
                \item adjusting for frame imperfections and nonresponse.
            \end{itemize}
        \item Information is available in the form of \alert{unit-level
            auxiliary variables} and known \alert{population totals} or
            \alert{group counts}.
        \item We distinguish the \alert{stage} at which the information is
            incorporated (role):
            \vspace{0.25em}
            \begin{itemize}
            \setlength{\itemsep}{3pt}
                \item \alert{Design stage} (sampler), e.g., information from
                    the sampling frame for identifying domains/strata
                \item \alert{Estimation stage} (analyst) $\Rightarrow$
                    post-stratification/calibration
            \end{itemize}
        \item Information is often used at \alert{both} stages.
    \end{itemize}
\end{frame}

% -----------------------------------------------------------------------------
\begin{frame}
    \begin{Bsp}[Ophthalmologists]

        \vspace{0.5em}

        Ophthalmologists by language region and gender (MedReg, 2022)

        \vspace{-0.5em}

        \begin{columns}
            \begin{column}{0.4\textwidth}
                \begin{center}
                \begin{tabular}{lr}
                    \toprule
                    \addlinespace
                    Language region & Count\\
                    \addlinespace
                    \midrule
                    German & 1'260\\
                    \addlinespace
                    French & 572\\
                    \addlinespace
                    Italian & 70\\
                    \bottomrule
                \end{tabular}
                \end{center}
            \end{column}
            \begin{column}{0.6\textwidth}
                \begin{center}
                \begin{tabular}{llr}
                    \toprule
                    \addlinespace
                    Language region & Gender & Count\\
                    \addlinespace
                    \midrule
                    German & Female & 416\\
                           & Male & 844\\
                    \addlinespace
                    French & Female & 189\\
                           & Male & 383\\
                    \addlinespace
                    Italian & Female & 26\\
                            & Male & 44\\
                    \bottomrule
                \end{tabular}
                \end{center}
            \end{column}
        \end{columns}
        \vspace{0.5em}
    \end{Bsp}
\end{frame}

% -----------------------------------------------------------------------------
\begin{frame}
    \begin{Def}[Grouping]\label{def:group}
        A \alert{grouping} $\mathscr{G}$ of the population $U$ is a partition
        \begin{equation*}
            \mathscr{G} = \{U_1,\ldots,U_G\},
            \qquad
            U = \bigcup_{g=1}^G U_g,
            \qquad
            U_g \cap U_{g'} = \varnothing
            \quad (g\neq g').
        \end{equation*}
    \end{Def}

    \vspace{0.25em}

    \begin{itemize}
    \setlength{\itemsep}{6pt}
        \item We use \alert{grouping} as a generic term. Depending on when the
            groups are defined, they may represent design strata,
            post-strata, or domains.
        \item Let $\mathscr H$ denote the stratification used at the design
            stage. We distinguish, for example:
            \vspace{0.25em}
            \begin{itemize}
            \setlength{\itemsep}{8pt}
                \item $\mathscr G = \mathscr H$: the groups coincide with the
                    design strata;
                \item $\mathscr G$ \alert{refines} $\mathscr H$: each group
                    $U_g$ is contained in one design stratum.
            \end{itemize}
    \end{itemize}
\end{frame}

\subsection{Post-stratification}

% -----------------------------------------------------------------------------
\begin{frame}
    \frametitle{\insertsectionnumber.\insertsubsectionnumber~\insertsubsection}
    \begin{itemize}
    \setlength{\itemsep}{10pt}
        \item Let $\mathscr G = \{U_1,\ldots,U_G\}$ denote a grouping of $U$
            with known group sizes $N_g$.
        \item The \alert{post-stratified estimator} of the population mean
            $\myu$ is
            \begin{equation*}
                \mypost
                =
                \sum_{g=1}^G W_g\,\mys_g,
                \qquad
                W_g=\frac{N_g}{N},
                \qquad
                \mys_g
                =
                \frac{1}{n_g}\sum_{i\in s_g}y_i.
            \end{equation*}
        \item If $\mathscr G=\mathscr H$, the post-strata coincide with the
            design strata and $\mypost=\mystrat$.
        \item Otherwise, the post-stratum sample sizes
            \begin{equation*}
                n_g(S) = \vert S\cap U_g \vert
            \end{equation*}
            are \alert{random variables}.
    \end{itemize}
\end{frame}

% -----------------------------------------------------------------------------
\begin{frame}
    \begin{itemize}
    \setlength{\itemsep}{9pt}
        \item In many applications, $\mathscr{G}$ refines the design
            stratification $\mathscr{H}$.
        \item Suppose the original sample is an \alert{\SRS} of fixed size $n$.
            \vspace{0.25em}
            \begin{itemize}
            \setlength{\itemsep}{6pt}
                \item The post-stratum sample sizes
                    \begin{equation*}
                        n_g(S)=|S\cap U_g|
                    \end{equation*}
                    are random, with
                    \begin{equation*}
                        \EE\{n_g(S)\}=nW_g.
                    \end{equation*}
                \item It is possible that $n_g(S)=0$ for some post-stratum.
                    This becomes unlikely when the expected post-stratum sample
                    sizes $nW_g$ are sufficiently large.
                \item In practice, post-strata with no or very few sampled
                    units may need to be \alert{combined}.
            \end{itemize}
    \end{itemize}
\end{frame}

% -----------------------------------------------------------------------------
\begin{frame}
    \begin{itemize}
    \setlength{\itemsep}{9pt}
        \item If the post-stratum sample sizes are sufficiently large, an
            \alert{approximate variance estimator} is
            \citep[][Remark 7.10.6]{sarndal_swensson_etal_1992}
            \begin{equation*}
                \varest\big(\mypost\big) \approx
                \underbrace{(1-f)\sum_{g=1}^G
                    W_g\,\frac{\widehat{\sigma}_g^2}{n}}_{
                    =\,\varest^{\mprop}(\mypost)} +
                \color{magenta}
                \underbrace{\color{black}
                    \frac{1-f}{n}
                    \sum_{g=1}^G
                    (1-W_g)\frac{\widehat{\sigma}_g^2}{n}}_{
                    \substack{\text{\color{magenta}additional contribution from}\\
                    \text{\color{magenta}random post-stratum sample sizes}}}
                \color{black}.
            \end{equation*}
            \begin{itemize}
                \item[\goto] The approximation is understood
                    \alert{conditional} on $n_g(S)>0$ for all $g = 1, \ldots,
                    G$.
            \end{itemize}
        \item The second term is of smaller order and becomes
            \alert{negligible} as the post-stratum sample sizes increase.
            Hence,
            \begin{equation*}
                \widehat{\var}\big(\mypost\big)
                \approx
                \widehat{\var}_{\mprop}\big(\mypost\big).
            \end{equation*}
    \end{itemize}
\end{frame}

% -----------------------------------------------------------------------------
\begin{frame}
    \begin{itemize}
    \setlength{\itemsep}{12pt}
        \item For sufficiently large post-stratum sample sizes, an \SRS with
            post-stratification is approximately as \alert{efficient} as
            stratified \SRS with \alert{proportional allocation}
            \citep{holt_smith_1979}.
        \item Stratification at the design stage is useful when we want to
            \vspace{0.25em}
            \begin{itemize}
            \setlength{\itemsep}{4pt}
                \item improve \alert{efficiency}, or
                \item control the sample sizes for important
                    \alert{domains}.
            \end{itemize}
        \item Post-stratification is more \alert{flexible}: different
            groupings can be used at the estimation stage for different
            variables or purposes.
        \item This flexibility is particularly useful in
            \alert{multipurpose surveys}.
    \end{itemize}
\end{frame}

\subsection{Calibration}

% -----------------------------------------------------------------------------
\begin{frame}
    \frametitle{\insertsectionnumber.\insertsubsectionnumber~\insertsubsection}
    \begin{itemize}
    \setlength{\itemsep}{8pt}
        \item Let $w_i=1/\pi_i$ denote the \alert{design weight} of sampled
            unit $i = 1, \ldots, n$.
        \item \alert{Calibration} replaces the design weights $w_i$ by
            \alert{calibration weights} $w_i^*$ that \citep{deville_sarndal_1992}
            \vspace{0.25em}
            \begin{enumerate}
            \setlength{\itemsep}{4pt}
                \item remain as close as possible to the design weights, and
                \item satisfy the \alert{calibration constraints}
                    \begin{equation*}
                        \sum_{i\in s}w_i^*\,\vx_i=\tvx, \qquad
                        \tvx=\sum_{i\in U}\vx_i,
                    \end{equation*}
                    where $\vx_i=(x_{i1},\ldots,x_{iq})^T$
                    is a vector of auxiliary variables.
            \end{enumerate}
        \item Thus, the weighted sample reproduces the known
            \alert{population totals} of the auxiliary variables.
    \end{itemize}
\end{frame}

% -----------------------------------------------------------------------------
\begin{frame}
    \begin{itemize}
    \setlength{\itemsep}{11pt}
        \item The \alert{calibration estimator} of the population total $\ty$
            is
            \begin{equation*}
                \widehat t_y^{\;cal} =
                \sum_{i\in s} w_i^* y_i.
            \end{equation*}
        \item Regarding the available \alert{auxiliary information}, we
            distinguish two cases:
            \vspace{0.5em}
            \begin{description}
            \setlength{\itemsep}{8pt}
                \item[Case 1)] $\vx_i$ is known for all $i\in U$; hence,
                    the population total $\tvx$ is also known.
                \item[Case 2)] $\vx_i$ is known only for $i\in s$, but the
                    population total $\tvx$ is known from an external source.
            \end{description}
        \item For calibration, the \alert{essential requirement} is that
            $\vx_i$ is available for the sampled units and that $\tvx$ is
            known.
    \end{itemize}
\end{frame}

%FIXME:

% -----------------------------------------------------------------------------
\begin{frame}
    \frametitle{\normalsize{Examples of choices for $\vx_i$} (and the
    corresponding population totals)}
    \begin{center}
    \begin{tabular}{lll}
        \toprule
        \addlinespace
        Case & Auxiliary vector ($i \in s$) & Known population total\\
        \addlinespace
        \midrule
        \addlinespace
        1 & $1$ & $N$\\
        \addlinespace
        2 & $x_i$ & $\tx$\\
        \addlinespace
        3 & $(1, \, x_i)^T$ & $(N, \, \tx)^T$\\
        \addlinespace
        4 & $\big(\indicator_i^D, \, \indicator_i^F, \, \indicator_i^I\big)^T$
            & $\big(N_D, N_F, N_I\big)^T$\\
        \addlinespace
        \bottomrule
    \end{tabular}
    \end{center}
    \vspace{0.5em}
    with indicator variables for language region: $\indicator_i^D=1$ if the
    unit is German-speaking and 0 otherwise; analogously for the French (F)
    and Italian (I) language regions. $N_D$, $N_F$, and $N_I$ denote the
    corresponding population sizes.
\end{frame}

% -----------------------------------------------------------------------------
\begin{frame}
    \begin{Bsp}[Language Regions]
        Calibration with
        \begin{equation*}
            \vx_i =
            \Big(1,\,\indicator_i^D,\,\indicator_i^F,\,
                \indicator_i^I\Big)^T
            \qquad \text{and} \qquad
            \tvx =
            \Big(N,\,N_D,\,N_F,\,N_I\Big)^T
        \end{equation*}
        leads to \alert{linearly dependent calibration variables}.
    \end{Bsp}
    \vspace{0.5em}
    \begin{itemize}
    \setlength{\itemsep}{10pt}
        \item Let $\bold X$ denote the $(n\times q)$ \alert{matrix} given by
            $\bold X = \big(\vx_1^T,\ldots,\vx_n^T\big)^T.$
            In this example, $\bold X$ does \alert{not have full column rank}.
        \item We can therefore either
            \vspace{0.25em}
            \begin{itemize}
            \setlength{\itemsep}{5pt}
                \item omit the constant $1$ from $\vx_i$ and $N$ from
                    $\tvx$, or
                \item omit one language-region indicator and the
                    corresponding population count.
            \end{itemize}
    \end{itemize}
\end{frame}

% -----------------------------------------------------------------------------
\begin{frame}
    \begin{itemize}
    \setlength{\itemsep}{10pt}
        \item The calibration constraints generally do not determine the
            weights $w_i^*$ uniquely. We therefore choose weights that remain
            \alert{as close as possible} to the design weights $w_i$.

        \item \citet{deville_sarndal_1992} consider several distance
            functions. One choice is the \alert{chi-square distance}
            \begin{equation*}
                D_{\chi^2}(\vw,\vw^*)
                =
                \sum_{i\in s}
                \frac{(w_i^*-w_i)^2}{q_iw_i},
                \qquad q_i>0.
            \end{equation*}

        \item The calibration weights are obtained by solving
            \begin{equation*}
                \min_{\vw^*}D_{\chi^2}(\vw,\vw^*)
                \qquad\text{subject to}\qquad
                \sum_{i\in s}w_i^*\vx_i=\tvx.
            \end{equation*}

        \item In what follows, we set $q_i\equiv1$.
    \end{itemize}
\end{frame}

% -----------------------------------------------------------------------------
\begin{frame}
    \begin{itemize}
    \setlength{\itemsep}{10pt}
        \item For the chi-square distance, the \alert{calibration weights}
            have the closed-form solution
            \begin{equation*}
                w_i^* =
                w_i\left[
                    1+
                    \Big(\tvx - \tvxs^{\,\mHT}\Big)^T
                    \left(
                        \sum_{j\in s}w_j\vx_j\vx_j^T
                    \right)^{-1}
                    \vx_i
                \right],
            \end{equation*}
            where
            \begin{equation*}
                \tvxs^{\,\mHT}=\sum_{i\in s}w_i\vx_i
            \end{equation*}
            is the Horvitz--Thompson estimator of $\tvx$.
        \item The resulting \alert{calibration estimator} of $\ty$ can be
            written as
            \begin{equation*}
                \widehat t_y^{\;cal} = \tysht_y +
                \big(\tvx - \tvxs\big)^T\vbs,
                \qquad \text{where} \qquad
                \vbs = \left(
                    \sum_{i\in s}w_i\vx_i\vx_i^T
                \right)^{-1}
                \sum_{i\in s}w_i\vx_i y_i.
            \end{equation*}
    \end{itemize}
\end{frame}

% -----------------------------------------------------------------------------
\begin{frame}
    \begin{itemize}
    \setlength{\itemsep}{11pt}
        \item The calibration estimator $\widehat t_y^{\;cal}$
            \vspace{0.25em}
            \begin{itemize}
            \setlength{\itemsep}{8pt}
                \item is generally \alert{not design-unbiased}, but is
                    design consistent under suitable regularity conditions;
                \item can be more \alert{efficient} than the HT estimator
                    when the auxiliary variables are informative about $y$;
                \item coincides, for the chi-square distance, with the
                    \alert{generalized regression estimator} (GREG).
            \end{itemize}
        \item \citet{deville_sarndal_1992} showed that calibration estimators
            based on a broad class of \alert{distance functions} are
            asymptotically equivalent.
        \item Variance estimation can be based on the residuals
            $y_i-\vx_i^T\vbs$.
    \end{itemize}
\end{frame}

% -----------------------------------------------------------------------------
\begin{frame}
    \frametitle{Why use calibration?}
    \begin{itemize}
    \setlength{\itemsep}{11pt}
        \item Calibration provides a systematic way to incorporate
            \alert{auxiliary information} into estimation.
        \item Important purposes include:
            \vspace{0.25em}
            \begin{itemize}
            \setlength{\itemsep}{6pt}
                \item improving \alert{efficiency} when the auxiliary
                    variables are informative about $y$;
                \item ensuring \alert{consistency with known population
                    totals};
                \item reducing the effects of \alert{nonresponse} or frame
                    imperfections.
            \end{itemize}
        \item For nonresponse adjustment, useful auxiliary variables should
            be related to the \alert{response propensity} and/or the
            \alert{variables of interest}; see
            \citet{sarndal_lundstrom_2005}.
    \end{itemize}
\end{frame}

% -----------------------------------------------------------------------------
\begin{frame}
    \frametitle{Historical context}
    \begin{itemize}
    \setlength{\itemsep}{10pt}
        \item The basic idea of adjusting survey weights to incorporate known
            population information predates \citet{deville_sarndal_1992}.
        \item An important early example is \alert{iterative proportional
            fitting} (raking), introduced by \citet{deming_stephan_1940} to
            adjust contingency tables to known marginal totals.
        \item \citet{deville_sarndal_1992} provided a general framework for
            \alert{calibration weighting}, based on calibration constraints
            and distance functions.
        \item \alert{Entropy balancing} \citep{hainmueller_2012}, widely used
            in the social sciences, can be viewed as a special case of
            calibration weighting using an entropy-type distance function.
    \end{itemize}
\end{frame}

% -----------------------------------------------------------------------------
\setcounter{section}{\the\numexpr\value{section}-1\relax}
\begin{frame}
    \frametitle{References}
    {\footnotesize
        \setstretch{1.1}
        \bibliographystyle{fhnw_en}
        \bibliography{master}
    }
\end{frame}

% -----------------------------------------------------------------------------
\lastpage

\end{document}
